\section*{Introducción}
La temperatura es una medida sobre la energia interna de un cuerpo, especificamente la energia cinetica de sus moleculas, la vibracion de estas. Entre mas energia cinetica la distancia promedio de sus moleculas aumenta. Esta distancia promedio se puede expresar como:
\[
	\Delta l/l_0,
\]
donde $\Delta l$ es la dilataci\'on del material y $l_0$ es su longitud original. Experimentalmente se ha encontrado que la distancia promedio de las moleculas es linealmente proporcional a su cambio de temperatura.
\begin{align*}
	\Delta l/l_0 \propto \Delta T \\
	\Delta l/l_0 = \alpha \Delta T 
\end{align*}
Este $\alpha$ es un factor dependiente del material y se le conoce como coeficiente promedio de expansion lineal.
En terminos de la diltacion la ecuacion es la siguiente:
\begin{equation}
	\label{eq:dilatacion}
	\Delta l = \alpha l_0 \Delta T
\end{equation}

En la practica se determina el coeficiente de expasion lineal para el acero y el cobre, ademas de compararlo con valores aceptados en la literatura.
